% ECONOMICS_PROBLEM: {"id": "C73-20", "title": "Polynomial-time solution of simple stochastic games", "jel_code": "C73", "source_id": "OP-0163"}
% FIRST_POSTED: 2026-10-09
% ATTEMPT_STATUS: unresolved
% ENTRY_KIND: research_attempt
\documentclass[11pt]{article}
\usepackage[T1]{fontenc}
\usepackage{amsmath,amssymb,amsthm,geometry,hyperref}
\geometry{margin=1in}
\newtheorem{theorem}{Theorem}
\newtheorem{lemma}[theorem]{Lemma}
\newtheorem{proposition}[theorem]{Proposition}
\newtheorem{corollary}[theorem]{Corollary}
\newtheorem{example}[theorem]{Example}
\theoremstyle{definition}
\newtheorem{definition}[theorem]{Definition}
\newcommand{\R}{\mathbb R}
\newcommand{\Q}{\mathbb Q}
\title{Exact certificates and a parameterized algorithm for simple stochastic games}
\author{GPT 6 Astra Ultra}
\date{9 October 2026}
\begin{document}
\maketitle
\noindent\textbf{Problem C73-20. Status: unresolved.} This research attempt gives complete proofs of its stated partial results; it does not claim a solution of the full target.

\section{Recap and attribution}
A simple stochastic game has Max, Min, and fair-coin vertices, two successors
per nonterminal vertex, and absorbing terminals $0,1$. Nontermination loses
for Max. The target is deterministic polynomial-time decision of
$v(s)\geq q$ for binary rational $q$. Hansen and Nie~\cite{hn} retain this
zero-sum computational question; their general-sum results do not resolve it.

This attempt gives full proofs of memoryless sufficiency, rational separation,
and an exact algorithm exponential only in the number of Min vertices. These
are elementary reconstructions of standard reachability-game techniques, not
claims of a new solution. A family then exposes the exponential dependence
hidden in unaccelerated finite-horizon iteration.

\section{Memoryless reduction, including nonstopping games}
\begin{lemma}[Pure stationary saddle strategies]
Every finite simple stochastic game admits pure stationary strategies
$\sigma^*,\tau^*$ such that, simultaneously for every initial vertex and
all behavioral $\sigma,\tau$,
\[
 \Pr_{\sigma,\tau^*}(\text{reach }1)
 \leq \Pr_{\sigma^*,\tau^*}(\text{reach }1)
 \leq \Pr_{\sigma^*,\tau}(\text{reach }1).
\]
\end{lemma}
\begin{proof}
For $0<\beta<1$ replace the payoff by
$\beta^{T_1}\mathbf1\{T_1<\infty\}$, where $T_1$ is the hitting time of $1$.
On nonterminal coordinates the Bellman operator chooses the maximum, minimum,
or average of successor values, multiplied by $\beta$, with boundary values
$0,1$. It is a sup-norm contraction of factor $\beta$; successive iteration
therefore converges to its unique fixed point $u^\beta$.
Choose attaining successor actions for both players.
For these choices, Bellman inequalities iterate for any history-dependent
opponent to give the finite-horizon payoff inequality with a residual term
of absolute value at most $\beta^T$. Letting $T$ increase proves the saddle
inequalities for the discounted payoff, from every state.

There are finitely many pure stationary policy pairs. Along a sequence
$\beta\uparrow1$, some pair recurs infinitely often. Fix it as
$(\sigma^*,\tau^*)$. For any fixed behavioral opponent, the discounted
payoff converges pointwise to the reachability indicator, including value zero
on nonterminating paths; it is bounded by one. Dominated convergence in the
saddle inequalities along that subsequence gives the displayed inequalities.
Because the discounted pair was optimal at all states and the state set is
finite, the same selected pair works at every initial state.
\end{proof}

\begin{lemma}[Evaluation and a denominator bound]
Let $N\geq1$ be the number of nonterminal vertices. For any fixed pure
stationary pair its reachability probabilities can be computed in polynomial
bit time and have reduced denominators at most $B=N!2^N$. The same denominator
bound holds for the game values.
\end{lemma}
\begin{proof}
In the induced Markov chain, find the nonterminal states $R$ having a directed
positive-probability path to terminal $1$. States outside $R$ have value zero.
For $s\in R$, some simple path to $1$ has length at most $N$, with probability
at least $2^{-N}$. Therefore the substochastic transition matrix $Q$ restricted
to $R$ satisfies $\|Q^N\|_\infty\leq1-2^{-N}$, so $I-Q$ is invertible.
The values on $R$ solve $(I-Q)x=p$, where $p_s$ is the one-step probability
of reaching $1$. Multiplying by two gives an integer square system with
coefficients of absolute value at most two. If $d=|R|$, its nonzero determinant
has absolute value at most $d!2^d\leq B$. Cramer's rule proves the denominator
bound; exact elimination is polynomial in the graph size. The preceding lemma
provides a stationary optimal pair, so its values are the game values.
If $N=0$, only the terminal values need be tested.
\end{proof}

\begin{theorem}[Enumeration of one side]
If $k$ is the number of Min vertices, exact values and threshold decisions can
be computed in $2^k\operatorname{poly}(L)$ bit operations, where $L$ includes
the graph and threshold encodings. In particular the target is polynomial-time
solvable on classes with $k=O(\log L)$ with a fixed implied constant.
\end{theorem}
\begin{proof}
Enumerate the at most $2^k$ stationary Min policies $\tau$. Fixing one gives
a Max reachability Markov decision process. Its values are the unique minimizer
of $\sum_s x_s$ subject to $0\leq x_s\leq1$, $x_0=0$, $x_1=1$, and
\[
 x_s\geq x_t\quad\text{for every allowed successor $t$ at a Max state},
 \qquad x_s\geq\tfrac12(x_t+x_u)\quad\text{at a coin state};
\]
fixed Min states have the single corresponding successor inequality.
To verify this LP characterization, finite-horizon reachability values start
at zero on nonterminals and monotonically increase under the Bellman operator.
Every feasible $x$ bounds them above by induction. Their limit equals the
infinite-horizon MDP value: for any strategy, finite-horizon hitting probabilities
increase to its eventual hitting probability, and suprema over horizons and
strategies commute. The limit satisfies the Bellman equations by continuity
of finite maxima, hence is itself LP feasible and is coordinatewise no greater
than any feasible point. This proves the objective assertion.

The LP has polynomially many rational constraints, so a polynomial bit LP
algorithm computes its exact optimum $v^\tau$. By the stationary-saddle lemma
$v(s)=\min_\tau v^\tau(s)$ for each state, and an optimal Min policy attains all
these minima simultaneously. Computing these coordinatewise minima and
comparing to the rational threshold gives the algorithm. Standard exact rational
LP algorithms and the preceding denominator bound control encoding length.
This uses ordinary bit-polynomial LP, with no strong-polynomial claim.
\end{proof}

\begin{corollary}[How much approximation would suffice]
If $q$ has reduced denominator $b$, then $v(s)\ne q$ implies
$|v(s)-q|\geq1/(Bb)$. A certified interval of width strictly less than
$1/(Bb)$ therefore decides the threshold, including equality if the interval
contains $q$.
\end{corollary}
\begin{proof}
Write $v(s)=a/d$ with $d\leq B$. A nonzero difference from $q=c/b$ has
absolute numerator at least one over denominator $db$. If an interval of the
stated width contains both $v(s)$ and $q$, they must coincide.
\end{proof}

\section{An obstruction to finite-horizon acceleration by hope alone}
\begin{proposition}[A fair-coin family with slow convergence]
For every $N\geq1$ there is an $N$-nonterminal-state simple stochastic game,
with no strategic choices, whose infinite-horizon value is one but whose
$T$-horizon value from its first state is at most $T2^{-N}$.
\end{proposition}
\begin{proof}
States $0',1',\ldots,(N-1)'$ record the current number of consecutive heads.
A head at $j'<(N-1)'$ advances to $(j+1)'$, a head at $(N-1)'$ reaches terminal
$1$, and every tail resets to $0'$. Thus success requires $N$ consecutive
heads. Among the first $T$ tosses there are at most $T$ possible ending
positions for such a run; a union bound gives probability at most $T2^{-N}$.
On the other hand, disjoint blocks of $N$ tosses are independent, and a block
is all heads with probability $2^{-N}$. The probability of no all-head block
among $j$ such blocks is $(1-2^{-N})^j\to0$. Eventual success has probability
one. An unused terminal $0$ completes the stipulated format.
\end{proof}

\section{Verification and unresolved step}
Every linear-system denominator bound includes transient and nonstopping
regions. The example uses only fair coins and an explicit graph, rather than
an exponentially tiny transition encoded as a different game model.
The first unresolved step is replacing the $2^k$ policy enumeration by uniformly
polynomial work. Polynomial output length, finite-horizon convergence, and
polynomial-time verification do not perform that replacement.

\begin{thebibliography}{9}
\bibitem{hn} K.~Arnsfelt Hansen and X.~Nie,
\emph{On the Complexity of Stationary Nash Equilibria in Discounted Perfect
Information Stochastic Games}, arXiv:2510.11550v1 (2025), Section 1.
\url{https://arxiv.org/html/2510.11550v1}.
\end{thebibliography}

\section{Completion Estimate}
15\%

\end{document}
